\section{Background}
\label{sec:background}

\subsection{Classical random walks and hitting times}
\label{sec:background-classical}

Let $G=(V,E)$ be a connected graph with $n=\lvert V\rvert$ vertices and let
$P$ be a row-stochastic transition matrix. A probability row vector $\pi$ is
stationary if $\pi P=\pi$. For a reversible chain, detailed balance holds:
$\pi_xP_{xy}=\pi_yP_{yx}$ for all $x,y$.

Given a non-empty marked set $M\subset V$, the hitting time is the expected
number of steps needed to reach $M$ from the stationary distribution,
\begin{equation}
  \HT(P,M)=\mathbb{E}_{X_0\sim\pi}\left[\min\{t\ge0:X_t\in M\}\right].
\end{equation}
For computations we use the equivalent conditioned form
\begin{equation}
  \HT(P,M)=\frac{1}{1-p_M}\sum_{x\notin M}\pi_x h_x,
  \qquad h_x=1+\sum_yP_{xy}h_y,\quad x\notin M,
  \label{eq:background-ht}
\end{equation}
with $p_M=\sum_{x\in M}\pi_x$. The linear system in \eqref{eq:background-ht}
is small for the graphs in this paper and can be solved directly. Standard
examples are $\HT=\Theta(N)$ on the complete graph, $\HT=\Theta(N^2)$ on the
cycle, and $\HT=\Theta(N\log N)$ on the two-dimensional grid; the hypercube has
$\HT=\Theta(N)$, and our data ($N=8$--$64$) give $\HT/N\approx1.15$--$1.34$.

\subsection{Szegedy walks and the detection--finding gap}
\label{sec:background-szegedy}

A Szegedy walk lifts a Markov chain to a unitary operator on two vertex
registers. With the convention used by AGJK,
\begin{equation}
  V(P)\lvert\bar0\rangle\lvert x\rangle
  =\sum_y\sqrt{P_{xy}}\lvert y,x\rangle,
  \qquad
  W(P)=V^\dagger(P)\,\mathrm{Shift}\,V(P)\,\mathrm{Ref},
  \label{eq:background-szegedy}
\end{equation}
where $\mathrm{Shift}\lvert x,y\rangle=\lvert y,x\rangle$ and
$\mathrm{Ref}=(2\lvert\bar0\rangle\langle\bar0\rvert-I)\otimes I$. The
discriminant matrix is
\begin{equation}
  D_{xy}=\sqrt{P_{xy}P_{yx}}
  =\langle\bar0|\langle x|W(P)|\bar0\rangle|y\rangle.
\end{equation}
For a reversible chain, $D=\mathrm{diag}(\pi)^{1/2}P\,\mathrm{diag}(\pi)^{-1/2}$,
so $D$ and $P$ share their spectrum and $D$ is symmetric.

Szegedy's search algorithm detects a marked vertex in
$O(\sqrt{\HT})$ steps. Detection is not the same as finding: measuring the
final state need not return a marked vertex. The gap is especially visible
when there are multiple marked vertices. This is the starting point of AGJK.

\subsection{The AGJK algorithm and extended hitting time}
\label{sec:background-agjk}

AGJK introduce the absorbing walk $P'$ and the interpolated walk
$P(s)=(1-s)P+sP'$, which stays at a marked vertex with probability $s$ instead
of moving. The corresponding quantum walk $W(s)$ uses an ancilla to record
whether the current vertex is marked. AGJK define the extended hitting time
$\HT^+$ through the eigenvalues and eigenvectors of the discriminant matrix
$D(s)$; if there is a single marked vertex, $\HT^+=\HT$, but with multiple
marked vertices $\HT^+$ can be much larger.

Two algorithms are given. The simpler one runs $W(s)$ for $O(\sqrt{\HT})$
steps and is supported by a numerical conjecture (Conjecture~3 in the paper):
for some $s$ and $t=O(\sqrt{\HT})$, measuring the vertex register returns a
marked vertex with constant probability. The proved algorithm combines the
interpolated walk with quantum fast-forwarding and achieves
$\widetilde O(\sqrt{\HT})$ time for any graph and any marked set
\citep{agjk2019}. We implement and test the first algorithm numerically; we
cite the second as the theoretical guarantee and do not re-prove it.

\subsection{Quantum cognition and quantum walks}
\label{sec:background-cognition}

Quantum cognition replaces classical probability with quantum probability in
models of judgment and decision. Pothos and Busemeyer review the program and
its explanations of order effects, conjunction and disjunction effects, and
contextuality \citep{pothos2022,busemeyer2024}. The models are usually
process-level: they do not assume that the brain implements a quantum circuit.

Quantum walks enter this literature in several ways. Chen, Ferro, and Sornette
use discrete-time quantum walks as evidence accumulators and discuss how
partial measurement produces decoherence \citep{chen2022}. Huang and Busemeyer
propose a quantum walk framework for multi-alternative decision making
\citep{huang2022}. Pavone and Faro use continuous-time quantum walks on
semantic graphs to model the emergence of insight through interference
\citep{pavone2025}. Epping et al.\ develop an open-system model of choice and
response time \citep{epping2023}, and Widdows et al.\ give quantum circuit
components for cognitive decision-making \citep{widdows2023}.

\begin{table}[t]
  \centering
  \small
  \begin{tabular}{@{}p{0.23\linewidth}p{0.25\linewidth}p{0.25\linewidth}p{0.19\linewidth}@{}}
    \toprule
    Paper & Model & Cognitive phenomenon & Role here \\
    \midrule
    AGJK 2019 & quantum walk search & marked-vertex search, $\HT^+$ & QIS core \\
    Chen et al.\ 2022 & DTQW & evidence accumulation, decoherence & measurement model \\
    Huang--Busemeyer 2022 & quantum walk & multi-alternative choice & decision mapping \\
    Pavone--Faro 2025 & CTQW on semantic graphs & insight, interference & semantic experiment \\
    Pothos--Busemeyer 2022 & quantum probability & context, order effects & motivation \\
    Epping et al.\ 2023 & open-system model & choice and response time & noise interpretation \\
    \bottomrule
  \end{tabular}
  \caption{Literature map. The technical core is AGJK; the other papers supply
  the cognitive interpretation and the measurement/noise models.}
  \label{tab:literature}
\end{table}

\subsection{The gap}
\label{sec:background-gap}

The two literatures meet at the level of metaphor but not at the level of
algorithmic structure. Quantum-walk search papers analyze worst-case
complexity and spectral quantities but do not ask what those quantities mean
for a cognitive process. Quantum-cognition papers use walks as evidence
accumulators but rarely exploit the marked-set structure, the extended hitting
time, or the difference between detection and finding. This paper treats the
marked-vertex search problem as a process model and uses its structural
quantities---$\HT$, $\HT^+$, the spectral gap, and the success curve---as the
bridge between the two literatures.
